% Standalone fragment: needs amsmath/amssymb, no project-specific macros.
% Suggested location: existing gain/loss appendix, after the Kraus and U(1)
% discussion. It has NOT been inserted into the manuscript.
\paragraph*{Gain and loss as Gaussian boundary limits.}
Although the conditional operators $\hat\chi^\dagger$ and $\hat\chi$ are
nilpotent, they arise as singular limits of invertible Gaussian operators.
For a normalized fermion mode, put
$\hat n=\hat\chi^\dagger\hat\chi$,
$\hat\Gamma=\hat\chi+\hat\chi^\dagger$, and
$\lambda=-\tfrac12\log\varepsilon$. Then
\begin{equation}
 \hat v_\pm(\varepsilon)
 =\hat\Gamma e^{\mp\lambda(2\hat n-\hat{\mathbf{1}})},\qquad
 \sqrt\varepsilon\,\hat v_+(\varepsilon)\to\hat\chi^\dagger,
 \quad
 \sqrt\varepsilon\,\hat v_-(\varepsilon)\to\hat\chi.
 \label{eq:GainLossGaussianBoundary}
\end{equation}
Each finite-$\varepsilon$ operator is a weak occupation filter followed
by a Majorana flip and induces a conjugation matrix in
$O(2N,\mathbb C)$ for $N$ complex modes. The endpoint is a projective
operator limit, not a finite singular element of that group. To obtain a
normalized instrument, use
$\hat M_0=[\hat{\mathbf{1}}-\hat n+\varepsilon\hat n]/\sqrt{1+\varepsilon^2}$
and
$\hat M_1=[\hat n+\varepsilon(\hat{\mathbf{1}}-\hat n)]/\sqrt{1+\varepsilon^2}$.
Feedback on the undesired outcome gives
$\{\hat\Gamma\hat M_0,\hat M_1\}\to\{\hat\chi^\dagger,\hat n\}$
for filling and
$\{\hat M_0,\hat\Gamma\hat M_1\}\to
\{\hat{\mathbf{1}}-\hat n,\hat\chi\}$ for depletion. For every fixed
finite record with nonzero limiting probability, the normalized
conditional state converges. This auxiliary regularization generally
breaks $U(1)$ at finite $\varepsilon$; it does not change the symmetry of
the exact protocol or justify exchanging the regulator and long-time
limits.
